% Part: incompleteness
% Chapter: incompleteness-provability
% Section: second-incompleteness-thm

\documentclass[../../../include/open-logic-section]{subfiles}

\begin{document}

\olfileid{inc}{inp}{2in}

\olsection{مبرهنة عدم الاكتمال الثانية}

نريد الآن ضبط قولنا إن $\Th{PA}$ لا تثبت اتساقها.
فعدم اتساق $\Th{PA}$ يكافئ $\Th{PA} \Proves \eq[{\OLMathNumeral{0}}][{\OLMathNumeral{1}}]$.
لذلك نعرّف عبارة الاتساق $\OCon[\Th{PA}]$ بأنها
الـ!!{sentence} $\lnot
\OProv[\Th{PA}](\gn{\eq[{\OLMathNumeral{0}}][{\OLMathNumeral{1}}]})$، فتأخذ دعوانا صورة المبرهنة
الآتية:

\begin{thm}
\ollabel{thm:second-incompleteness}
متى كانت $\Th{PA}$ متسقة، امتنع على $\Th{PA}$ !!{derive}
الصيغة $\OCon[\Th{PA}]$.
\end{thm}

ولا بد من التنبه إلى اعتماد هذه المبرهنة على التمثيل المختار
لـ$\OCon[\Th{PA}]$، أي على كيفية تمثيل $\OProv[\Th{PA}]({\OLId{latin-ordinary}{y}})$.
ولن نحتاج من $\OProv[\Th{PA}]({\OLId{latin-ordinary}{y}})$ إلا استيفاء شروط
!!{derivability} الثلاثة. فكل نظرية لها محمول
!!a{derivability} يستوفيها يجري عليها البرهان نفسه.

ويكشف مخطط غودل للحجة كيف تأتي النتيجة عقب مقدماتها مباشرة.
لتكن $!G_\Th{PA}$ جملة غودل المبنية في برهان
\olref[1in]{thm:first-incompleteness}. قد أثبتنا: «إذا كانت
$\Th{PA}$ متسقة، امتنعت على $\Th{PA}$ !!{derive} الصيغة
$!G_\Th{PA}$». فلو صُوّر هذا البرهان \emph{داخل} $\Th{PA}$
لحصلنا على برهان
\[
\OCon[\Th{PA}] \lif \lnot \OProv[\Th{PA}](\gn{!G_\Th{PA}}).
\]
فلنفرض أن النظام $\Th{PA}$ !!{derive}s $\OCon[\Th{PA}]$.
حينئذ !!{derive}s أيضًا $\lnot \OProv[\Th{PA}](\gn{!G_\Th{PA}})$.
وبما أن $!G_\Th{PA}$ جملة غودل، فهذه مكافئة لـ$!G_\Th{PA}$؛
إذن !!{derive}s النظام $\Th{PA}$ الصيغة $!G_\Th{PA}$.

لكن اتساق $\Th{PA}$ يمنعها، كما علمنا، من !!{derive}
الصيغة $!G_\Th{PA}$!{} فمتى كانت $\Th{PA}$ متسقة امتنع
عليها !!{derive} عبارة اتساقها $\OCon[\Th{PA}]$.

ولنضبط الحجة نأخذ $!G_\Th{PA}$ جملة غودل للنظرية $\Th{PA}$،
ونستعمل شروط !!{derivability} (P1)--(P3) لإثبات أن $\Th{PA}$
!!{derive}s $\OCon[\Th{PA}] \lif !G_\Th{PA}$.
ومن ذلك يتبيّن أن $\Th{PA}$ لا يمكنها !!{derive} $\OCon[\Th{PA}]$.
وهذا ترتيب البرهان في $\Th{PA}$، مع حذف المؤشرات $\Th{PA}$
السفلية طلبًا للاختصار:
\begin{align}
& !G \liff \lnot \OProv(\gn{!G}) \ollabel{G2-1}\\
& \qquad\text{إن~$!G$ جملة غودل}\notag \\
& !G \lif \lnot \OProv(\gn{!G}) \ollabel{G2-2}\\
  & \qquad\text{من \olref{G2-1}} \notag\\
& !G \lif
  (\OProv(\gn{!G}) \lif \lfalse) \ollabel{G2-3}\\
  & \qquad\text{من \olref{G2-2} بالمنطق}\notag\\
& \OProv(\gn{
    !G \lif
    (\OProv(\gn{!G}) \lif \lfalse)
  }) \ollabel{G2-4}\\
  & \qquad\text{من \olref{G2-3} بحسب الشرط P1} \notag\\
& \OProv(\gn{!G}) \lif
  \OProv(\gn{
    (\OProv(\gn{!G}) \lif \lfalse)
    }) \ollabel{G2-5}\\
  & \qquad\text{من \olref{G2-4} بحسب الشرط P2} \notag\\
& \OProv(\gn{!G}) \lif (\OProv(\gn{\OProv(\gn{!G})}) \lif \OProv(\gn{\lfalse})) \ollabel{G2-6}\\
  & \qquad\text{من \olref{G2-5} بحسب الشرط P2 والمنطق} \notag\\
& \OProv(\gn{!G}) \lif
  \OProv(\gn{\OProv(\gn{!G})}) \ollabel{G2-7}\\
   & \qquad\text{بحسب P3} \notag\\
& \OProv(\gn{!G}) \lif \OProv(\gn{\lfalse}) \ollabel{G2-8}\\
  & \qquad \text{من \olref{G2-6} و\olref{G2-7} بالمنطق}\notag\\
& \OCon \lif \lnot \OProv(\gn{!G}) \ollabel{G2-9}\\
  & \qquad\text{بعكس نقيض \olref{G2-8} و$\OCon \ident \lnot \OProv(\gn{\lfalse})$}\notag \\
& \OCon \lif !G \notag\\
  & \qquad\text{من \olref{G2-1} و\olref{G2-9} بالمنطق}\notag
\end{align}
والمقصود بالمنطق في هذه الخطوات حقائق أولية من المنطق القضوي.
فالخطوة \olref{G2-3} تعتمد
$\Proves \lnot!A \liff (!A\lif \lfalse)$،
والخطوة \olref{G2-8} تعتمد
$!A \lif (!B \lif !C), !A \lif !B \Proves !A \lif !C$.
وأما تطبيق P2 في \olref{G2-5} و\olref{G2-6} فبحالتين من
$\OProv(\gn{!A \lif !B}) \lif
(\OProv(\gn{!A}) \lif \OProv(\gn{!B}))$:
نأخذ في الأولى $!A \ident !G$ و$!B \ident \OProv(\gn{!G}) \lif \lfalse$،
وفي الثانية $!A \ident \OProv(\gn{!G})$ و$!B \ident \lfalse$.

وعلى وجه أعم نصوغ مبرهنة عدم الاكتمال الثانية هكذا:

\begin{thm}
\ollabel{thm:second-incompleteness-gen} إذا كانت $\Th{T}$
نظرية متسقة و!!{axiomatized} تمتد $\Th{Q}$، وكانت
$\OProv[\Th{T}]({\OLId{latin-ordinary}{y}})$ صيغة تستوفي شروط !!{derivability}
P1--P3 للنظرية $\Th{T}$، امتنع !!{derive} النظام
$\Th{T}$ للصيغة $\OCon[\Th{T}]$.
\end{thm}

\begin{prob}
أثبت أن النظام $\Th{PA}$ !!{derive}s الصيغة
$!G_{\Th{PA}} \lif \OCon[\Th{PA}]$.
\end{prob}

\begin{digress}
فحاصل الأمر أن النظرية الرياضية المتسقة «المعقولة» لا يمكنها
!!{derive} عبارة اتساقها. لنأخذ $\Th{T}$ نظرية تشمل
$\Th{Q}$ وتشمل استدلال هيلبرت «التناهي»، أيًّا كان تحديده.
إذا عجزت $\Th{T}$ بجملتها عن !!{derive} عبارة اتساق $\Th{T}$،
فجزؤها التناهي أولى بالعجز عن !!{derive} عبارة اتساق $\Th{T}$.
وبهذا المعنى يمتنع البرهان التناهي على اتساق «الرياضيات كلها».

ومع ذلك، يحتمل لفظ «التناهي» أكثر من تفسير. لذلك أجاز غودل
في بحثه سنة \OLClassicalDigits{1931} أن يكون بعض ما نعدّه تناهيًا خارج الرياضيات التي
طلب هيلبرت صياغتها صوريًا. وكان في هذا متسامحًا؛ إذ يعسر اليوم
تصور استدلال يصح عدّه تناهيًا ولا يقبل الصياغة الصورية في
$\Th{ZFC}$، مثلًا، وهي نظرية مجموعات زرميلو--فرينكل مع بديهية
الاختيار.
\end{digress}

\end{document}
